\chapter{Exercicis proposats}
\label{cap:exercicis}

\begin{enumerate}
\item Reformuleu els enunciats següents fent ús de variables i constants:

\begin{enumerate}
\item El cub d'un nombre senar és senar.

\item Si el quadrat d'un nombre enter no és divisible per 4, aleshores
aquest nombre és senar.

\item El producte de dos nombres consecutius és parell.
\end{enumerate}

\item Quines de les expressions següents són proposicions i quines predicats?

\begin{enumerate}
\item Si $x\geq2$, llavors $x^{3}\geq1.$

\item $(x+y)^{2}=x^{2}+2xy+y^{2}$.

\item Si $w=3$, llavors $z^{w}\neq0$.
\end{enumerate}

\item Quins dels següents enunciats són predicats i quins no ho
són? Justifica la teva resposta.\qquad

\begin{enumerate}
\item $x$ és divisible per 3.

\item La suma dels nombres \thinspace x\thinspace\ i \thinspace2.

\item $x^{2}-y^{2}$.

\item $x+2<y-3$.
\end{enumerate}

\item Considereu els enunciats següents: $p=$ \enquote{ Estic
content}, $q=$ \enquote{ Estic veient una
pel·lícula}\ i $r=$ \enquote{
Estic menjant espaguetis}. Expresseu en paraules els
enunciats següents: (a) $r\longrightarrow p$; (b) $p\longleftrightarrow
q$; (3) $q\vee\left(  r\longrightarrow p\right)  $; (4) $\left(
q\longrightarrow\lnot p\right)  \wedge\left(  r\longrightarrow\lnot p\right)
$.

\item Considereu els enunciats següents: $p=$ \enquote{ L'Eduard
té els cabells vermells}, $q=$ \enquote{ L'Eduard
té un nas gran}\ i $r=$ \enquote{ A l'Eduard li
agrada menjar crispetes}. Tradueix els enunciats següents
a símbols: (a) \enquote{ A l'Eduard no li agrada menjar
crispetes}; (b) \enquote{ L'Eduard té un nas gran
i els cabells vermells, o bé té un nas gran i li agrada menjar
crispetes}.

\item Construeix les taules de veritat de les següents expressions
lògiques: (a) $p\longrightarrow(q\wedge p)$; (b) $(p\wedge
q)\longleftrightarrow(p\vee\lnot r)$.

\item Considerem els predicats $P(x,y)=$ \enquote{$x+y=4$%
}\ i $Q(x,y)=$ \enquote{$x<y$}.
Troba els valors de veritat de les següents expressions lògiques: (a)
$P(x,y)\wedge Q(x,y)$ i (b) $P(x,y)\longrightarrow\lnot Q(x,y)$ quan usem els
valors següents (1) $x=3,y=1$; (2) $x=1,y=2$.

\item Proveu les implicacions següents:

\begin{enumerate}
\item $(A\longrightarrow B)\wedge A$ $\Longrightarrow B$

\item (i) $A\wedge B$ $\Longrightarrow A$; (ii) $A\wedge B$ $\Longrightarrow
B$

\item (i) $A\Longrightarrow A\vee B$; (ii) $B\Longrightarrow A\vee B$

\item (i) $\left(  A\vee B\right)  \wedge\lnot B\Longrightarrow A$; (ii)
$\left(  A\vee B\right)  \wedge\lnot A\Longrightarrow B$

\item (i) $A\longleftrightarrow B\Longrightarrow A\longrightarrow B$; (ii)
$A\longleftrightarrow B\Longrightarrow B\longrightarrow A$

\item $\left(  A\longrightarrow B\right)  \wedge\left(  B\longrightarrow
A\right)  \Longrightarrow A\longleftrightarrow B$
\end{enumerate}

\item Proveu les equivalències lògiques següents:

\begin{enumerate}
\item $\lnot\left(  \lnot A\right)  \Longleftrightarrow A$ (llei de la doble negació)

\item $A\wedge B\Longleftrightarrow B\wedge A$ (llei commutativa de $\wedge$)

\item $A\vee B\Longleftrightarrow B\vee A$ (llei commutativa de $\vee$)

\item $\left(  A\wedge B\right)  \wedge C\Longleftrightarrow A\wedge\left(
B\wedge C\right)  $ (llei associativa de $\wedge$)

\item $\left(  A\vee B\right)  \vee C\Longleftrightarrow A\vee\left(  B\vee
C\right)  $ (llei associativa de $\vee$)

\item $A\vee\left(  B\wedge C\right)  \Longleftrightarrow\left(  A\vee
B\right)  \wedge\left(  A\vee C\right)  $ (llei distributiva de $\vee$
respecte de $\wedge$)

\item $A\wedge\left(  B\vee C\right)  \Longleftrightarrow\left(  A\wedge
B\right)  \vee\left(  A\wedge C\right)  $ (llei distributiva de $\wedge$
respecte de $\vee$)

\item $A\longrightarrow B\Longleftrightarrow\lnot B\longrightarrow\lnot A$
(llei del contrarecíproc)

\item $A\longleftrightarrow B\Longleftrightarrow\left(  A\longrightarrow
B\right)  \wedge\left(  B\longrightarrow A\right)  $ (llei del bicondicional)

\item $\left(  A\vee B\right)  \longrightarrow C\Longleftrightarrow\left(
A\longrightarrow C\right)  \wedge\left(  B\longrightarrow C\right)  $

\item $\lnot(A\wedge B)\Longleftrightarrow\lnot A\vee\lnot B$ (llei de De Morgan)

\item $\lnot(A\vee B)\Longleftrightarrow\lnot A\wedge\lnot B$ (llei de De Morgan)
\end{enumerate}

\item Simplifiqueu les afirmacions següents. Podeu fer ús de les
equivalències de l'exercici 9 a més de les equivalències
comentades a teoria i pràctica: (a) $\lnot(X\longrightarrow\lnot Y)$; (b)
$\left(  Y\wedge Z\right)  \longrightarrow Y$; (c) $\lnot(X\longrightarrow
Y)\vee Y$; (d) $\left(  X\longrightarrow Y\right)  \vee Y$.

\item Escriu la negació dels enunciats següents: (a) $e^{3}>0$; (b)
$\sin\left(  \frac{\pi}{2}\right)  >0$ i $\tan0\geq0$; (c) $y-3>0$ implica
$y^{2}+9>6y$.

\item Per a cadascun dels arguments següents, si és vàlid, doneu
una deducció, i si no és vàlid, mostreu el perquè.

\begin{enumerate}
\item $%
\begin{tabular}
[c]{ll}%
$P_{1}:$ & $p\longrightarrow q$\\
$P_{2}:$ & $\lnot r\longrightarrow\lnot q$\\
$P_{3}:$ & $s\longrightarrow t$\\
$P_{4}:$ & $p\vee s$\\\hline
$C:$ & $r\vee t$%
\end{tabular}
\ \ \ \ \ $

\item $%
\begin{tabular}
[c]{ll}%
$P_{1}:$ & $p\longrightarrow q$\\
$P_{2}:$ & $\lnot r\longrightarrow\left(  s\longrightarrow t\right)  $\\
$P_{3}:$ & $r\vee\left(  p\vee t\right)  $\\
$P_{4}:$ & $\lnot r$\\\hline
$C:$ & $q\vee s$%
\end{tabular}
\ \ \ \ \ \ $

\item $%
\begin{tabular}
[c]{ll}%
$P_{1}:$ & $\lnot p\longrightarrow\left(  q\longrightarrow\lnot r\right)  $\\
$P_{2}:$ & $r\longrightarrow\lnot p$\\
$P_{3}:$ & $\left(  \lnot s\vee p\right)  \longrightarrow\lnot\lnot r$\\
$P_{4}:$ & $\lnot s$\\\hline
$C:$ & $\lnot q$%
\end{tabular}
\ \ \ \ \ $
\end{enumerate}

\item Per a cadascun dels arguments següents, si és vàlid, doneu
una deducció, i si no és vàlid, mostreu el perquè.

\begin{enumerate}
\item Si el rellotge està avançat, llavors en Joan va arribar abans de
les deu i va veure sortir el cotxe de l'Andreu. Si l'Andreu no diu la veritat,
llavors en Joan no va veure sortir el cotxe de l'Andreu. L'Andreu diu la
veritat o estava en l'edifici en el moment del crim. El rellotge està
avançat. Per tant, l'Andreu estava en l'edifici en el moment del crim.

\item Si $\alpha$ i $\beta$ són dos angles iguals, llavors $\alpha=45%
%TCIMACRO{\U{ba}}%
%BeginExpansion
{{}^o}%
%EndExpansion
$. Si $\beta=45%
%TCIMACRO{\U{ba}}%
%BeginExpansion
{{}^o}%
%EndExpansion
$, llavors $\alpha=90%
%TCIMACRO{\U{ba}}%
%BeginExpansion
{{}^o}%
%EndExpansion
$. O $\beta$ és recte o bé $\beta=45%
%TCIMACRO{\U{ba}}%
%BeginExpansion
{{}^o}%
%EndExpansion
$. $\beta$ no és recte. Per tant, $\alpha$ i $\beta$ no són iguals i
cap d'ells és recte.
\end{enumerate}

\item Escriviu una deducció per a cadascun dels arguments següents,
tots ells són vàlids. Indiqueu si les premisses són consistents o inconsistents:

\begin{enumerate}
\item No és el cas que la roba sigui molesta o no sigui barata. La roba no
és barata o no està de moda. Si la roba no està de moda, és
una ximpleria. Per tant, la roba és una ximpleria.

\item Si al Marc li agrada la pizza, li agrada la cervesa. Si al Marc li
agrada la cervesa, no li agrada l'arengada. Si al Marc li agrada la pizza, li
agraden les arengades. Al Marc li agrada la pizza. Per tant, li agrada la
pizza d'arengada.
\end{enumerate}

\item Trobeu la fal·làcia o fal·
làcies en cadascun dels arguments següents:

\begin{enumerate}
\item Si la meva tortuga menja una hamburguesa, es posa malalta. Si la meva
tortuga es posa malalta, llavors no és feliç. Per tant, la meva
tortuga es posa malalta.

\item Si Frida es menja una granota, Susana menjarà una serp. Frida no
menja una granota. Per tant, Susana no menja una serp.
\end{enumerate}

\item Suposem que els valors possibles de $x$ són totes les persones.
Considerem els predicats següents: $Y(x)=$ \enquote{$x$ té el cabell verd},
$Z(x)=$ \enquote{$x$ li agrada les croquetes} i $W(x)=$ \enquote{$x$ té una granota de
mascota}, $L(x,y)=$\enquote{$x$ és tan ràpid com $y$}, $M(x,y)=$\enquote{$x$ és
tan car com $y$} i $N(x,y)=$\enquote{$x$ és tan vell com $y$}. Tradueix les
expressions següents en paraules:%
\[%
\begin{tabular}
[c]{llllll}%
(1) & $\left(  \forall x\right)  Y(x)$ &  &  & (3) & $\left(  \exists
x\right)  \left(  Y(x)\longrightarrow Z(x)\right)  $\\
(2) & $\left(  \exists x\right)  Z(x)$ &  &  & (4) & $\left(  \forall
x\right)  \left(  W(x)\longleftrightarrow\lnot Z(x)\right)  $%
\end{tabular}
\ \ \ \
\]


\item Suposem que els valors possibles de $x$ i $y$ són tots els vehicles.
Considerem els predicats següents: $L(x,y)=$\enquote{$x$ és tan ràpid com
$y$}, $M(x,y)=$\enquote{$x$ és tan car com $y$} i $N(x,y)=$\enquote{$x$ és tan vell
com $y$}. Tradueix les expressions següents en paraules:%
\[%
\begin{tabular}
[c]{llllll}%
(1) & $\left(  \exists x\right)  \left(  \forall y\right)  L(x,y)$ &  &  &
(3) & $\left(  \exists y\right)  \left(  \forall x\right)  \left(  L(x,y)\vee
N(x,y)\right)  $\\
(2) & $\left(  \forall x\right)  \left(  \exists y\right)  M(x,y)$ &  &  &
(4) & $\left(  \forall y\right)  \left(  \exists x\right)  \left(  \lnot
M(x,y)\longrightarrow L(x,y)\right)  $%
\end{tabular}
\ \ \
\]


\item Expresseu formalment els enunciats següents:%
\[%
\begin{tabular}
[c]{ll}%
(1) & La gent és simpàtica.\\
(2) & A ningú li agraden els gelats i els embutits.\\
(3) & Em va agradar un dels llibres que vaig llegir l'estiu passat.\\
(4) & Existeix una vaca tal que, si té quatre anys, no té taques
blanques.\\
(5) & Per a cada fruita, hi ha una fruita més madura que ella.
\end{tabular}
\ \ \
\]


\item Escriu una negació de cada enunciat. No escriviu la paraula
\enquote{ no}\ davant de qualsevol dels objectes que
es quantifiquen (per exemple, no escriviu \enquote{ No tots els nois
són bons}\ per a la part (1) d'aquest exercici).%
\[%
\begin{tabular}
[c]{ll}%
(1) & Tots els nois són bons.\\
(2) & Hi ha ratpenats que pesen 3 kg o més.\\
(3) & L'equació $x^{2}-2x>0$ val per a tots els nombres reals $x$.\\
(4) & Hi ha un nombre enter $n$ tal que $n^{2}$ és un nombre perfecte.\\
(5) & Cada casa té una porta que és blanca.
\end{tabular}
\ \ \ \
\]


\item Assumint com a domini de les variables $x$ i $y$ el conjunt $U$,
escriviu una deducció per a cadascun dels arguments següents:

\begin{enumerate}
\item
\begin{tabular}
[c]{ll}%
$P_{1}:$ & $\left(  \forall x\right)  \left(  R(x)\longrightarrow C(x)\right)
$\\
$P_{2}:$ & $\left(  \forall x\right)  \left(  T(x)\longrightarrow R(x)\right)
$\\\hline
$C:$ & $\left(  \forall x\right)  \left(  \lnot C(x)\longrightarrow\lnot
T(x)\right)  $%
\end{tabular}
$\ \ \ \ \ $

\item
\begin{tabular}
[c]{ll}%
$P_{1}:$ & $\left(  \forall x\right)  \left(  \exists y\right)  \left(
E(x)\longrightarrow\left(  M(x)\vee N(x)\right)  \right)  $\\
$P_{2}:$ & $\lnot\left(  \forall x\right)  \ M(x)$\\
$P_{3}:$ & $\left(  \forall x\right)  \ E(x)$\\\hline
$C:$ & $\left(  \exists x\right)  \ N(x)$%
\end{tabular}

\end{enumerate}

\item Escriviu una deducció per a cadascun dels arguments següents:

\begin{enumerate}
\item Tots els peixos amb moltes espines no són agradables de menjar. Tots
els peixos amb moltes espines són peixos blancs. Per tant, tots els peixos
que siguin agradables de menjar no tenen moltes espines.

\item Tots els estudiants d'un institut de secundària que assisteixen a
classes d'ampliació són genials. Hi ha un estudiant de l'institut que
és intel·ligent i no és genial. Per tant, hi ha un
estudiant de l'institut que és intel·ligent i no
assisteix a classes d'ampliació.
\end{enumerate}

\item Construïu definicions de \enquote{quadrat perfecte},
\enquote{paral·lela} i \enquote{funció real de variable real}. Quins
són els dominis de les variables, quins termes s'han d'assumir com a coneguts?

\item Definiu la igualtat de dos nombres racionals, de dues rectes en el pla i
de dues funcions reals de variable real.

\item Considerem dos nombres enters $x$ i $y$, llavors proveu directament que

\begin{enumerate}
\item si $x$ i $y$ són parells, aleshores $x+y$ és parell.

\item si $x$ i $y$ són senars, aleshores $x+y$ és parell.

\item si un d'ells és senar i l'altre és parell, aleshores $x+y$
és senar.
\end{enumerate}

\item Si $n$ és un nombre natural imparell, proveu directament que $n^{2}
$ és de la forma $8k+1$, per a algun enter $k\geq0$.

\item Demostreu directament que les solucions reals de l'equació $ax^{2}+bx+c=0
$, on $a,b,c\in\mathbb{R}$, $a\neq0$ i $b^{2}-4ac\geq0$, són expressades per la
fórmula
\[
x=\frac{-b\pm\sqrt{b^{2}-4ac}}{2a}.
\]


\item Proveu cap a enrere:

\begin{enumerate}
\item Que qualsevol enter divideix zero.

\item Que per a tots $x,y\in\mathbb{R}$, $xy=0$ si i només si $x=0$ o
$y=0$.
\end{enumerate}

\item Proveu per contrarecíproc:

\begin{enumerate}
\item Que per a qualssevol nombres reals $a$ i $b$, $\left\vert a\right\vert
<\left\vert b\right\vert $ si i només si $a^{2}<b^{2}$.

\item Que per a qualssevol nombres reals $x$ i $y$, $x^{2}+y^{2}=0$ si i
només si $x=y=0$.
\end{enumerate}

\item Considerem que $a$ és un nombre racional i $b$ és irracional.
Llavors, proveu per reducció a l'absurd:

\begin{enumerate}
\item Que$~a+b$ és irracional

\item Si $a\neq0$, aleshores $ab$ és irracional.
\end{enumerate}

\item Proveu per reducció a l'absurd que hi ha infinits nombres primers.

\item Donats $a,b,c\in\mathbb{Z}$, proveu per reducció a l'absurd que si
$a$ no divideix $bc$, llavors $a$ no divideix $b$.

\item És pot demostrar que $n^{3}-n$ és múltiple de 3, $n^{5}-n$
és múltiple de $5$ i $n^{7}-n$ és múltiple de 7, però
$n^{k}-n$ és múltiple de $k$?

\item Proveu per inducció:

\begin{enumerate}
\item Si $n\in\mathbb{N}$, llavors $1+3+5+\cdots+\left(  2n-1\right)  =n^{2}$.

\item Si $n\in\mathbb{N}$ i $n\geq5$, $4n>n^{4}$.

\item Si $n$ és un nombre enter no negatiu, llavors $5\mid\left(
n^{5}-n\right)  $.
\end{enumerate}

\item (Desigualtat triangular) Si $x,y\in\mathbb{R}$, proveu que $\left\vert
x+y\right\vert \leq\left\vert x\right\vert +\left\vert y\right\vert $.

\item Proveu que no existeixen dos nombres enters $m$ i $n$ tals que
$14m+21n=100$.

\item Considerem que $a,b\in\mathbb{N}$. Llavors hi ha un únic
$d\in\mathbb{N}$ tal que: Un enter $m$ és múltiple de $d$ si i
només si $m=ax+by$ per alguns $x,y\in\mathbb{Z}$.

\item Definim per recursió per a qualsevol enter no negatiu $k\,:$

\begin{itemize}
\item $u_{0}=0$,

\item $u_{k+1}=3u_{k}+3^{k}$.
\end{itemize}

Proveu que $u_{n}=n\cdot3^{n-1}$ per a tot enter no negatiu $n$.

\item Demostreu que el principi de bona ordenació implica el principi d'inducció.

\begin{description}
\item \textbf{Suggeriment}: Considereu una proposició $p(n)$ sobre el
nombre natural $n$. Suposeu que $p(1)$ és vertadera, i també que per a
tots els $n$, $p(n)\longrightarrow p(n+1)$. Heu de demostrar que per a tot
$n$, $p(n)$ és vertadera. Ara feu la prova per reducció a l'absurd.
Suposeu que hi ha algun $n$ pel qual $p(n)$ és falsa. Llavors, considereu
el conjunt $K$ que té com a elements els nombres naturals $n$ que no
compleixen $p(n)$. Ara heu de concloure que això és una contradicció.
\end{description}
\end{enumerate}